%#####################################################
\section{Introduction}
\label{sec:intro}
%#####################################################

With the current and upcoming generations of large-scale optical imaging surveys, such as the Hyper Suprime-Cam Subaru Strategic Program (HSC-SSP, \cite{HSC2022datarelease3}), the Dark Energy Survey (DES, \cite{DES2018datarelease1}), the Kilo-Degree Survey (KiDS, \cite{KiDS2025datarelease5}), the Nancy Grace Roman Space Telescope \cite{roman2019wfirst}, Euclid \cite{EuclidI2024overview} and the Vera C. Rubin Observatory's Legacy Survey of Space and Time (LSST, \cite{LSST2019survey}) probing hundreds of millions to billions of galaxies with accurate shape information, the strength of cosmological parameter inference from weak lensing analysis is becoming increasingly important. Some of the most precise single measurements of the lensing amplitude $S_8\equiv\sigma_8\sqrt{\Omega_m/0.3}$ currently come from cosmic shear analyses \cite{CosmoShearLi2023HSCY3, Dalal2023CosmoShearHSC, Amon2022CosmoShearDES, Secco2022CosmoShearDES, Asgari2021KidsCosmoShear, wright2025kidslegacyCosmoShear}, where one measures the correlation in the distortion of galaxy images due to weak gravitational lensing. 

These optical surveys rely on characterization of their normalized sample redshift distributions projected along the line of sight $n(z)$ when measuring two-point statistics of galaxy clustering and gravitational shear fields for cosmological parameter inference, most often performed through the two-point correlation functions \cite{CosmoShearLi2023HSCY3} or the angular power spectrum \cite{Dalal2023CosmoShearHSC}. Since optical surveys rely on broadband photometry covering optical and near-infrared wavelengths, narrow spectral galaxy features such as emission lines are difficult to identify due to the limited coverage per observed object and the broadband nature of the optical filters. Therefore, imaging surveys employ photometric redshift (photo-$z$) algorithms to infer individual redshifts. Typically, photometric redshifts are used to split the galaxies into tomographic redshift bins, and the true redshift distribution along the line of sight needs to be estimated for each bin. 
One can reconstruct the overall redshift distribution $n(z)$ from observed data for each tomographic bin. In the case of the fiducial HSC Year 3 analysis \cite{HSCClusteringRau2023, CosmoShearLi2023HSCY3, Dalal2023CosmoShearHSC}, four tomographic bins are used, spanning linear bins from $z_p\sim0.3$ to $z_p\sim1.5$, where $z_p$ here is a photometric redshift from a given photo-$z$ method. Henceforth, Bin 1 refers to $0.3\leq z_p<0.6$, Bin 2 to $0.6\leq z_p<0.9$, Bin 3 to $0.9\leq z_p<1.2$ and Bin 4 to $1.2\leq z_p<1.5$ respectively. 
Many different algorithms have been developed for photometric redshifts, ranging from inferring redshifts based on prior knowledge of galaxy spectral energy distributions (SEDs) (\cite{LePhare1, LePhare2, Brammer2008EAZY, MizukiHSCTanaka2015}) to mapping redshifts to photometry information with machine learning (\cite{CarrascoKind2014PhotozMLz, Hsieh2014DEMp, Collister2004ANNz}). Often, since the performance of the algorithms does not allow for perfect point estimates of the redshift, photo-$z$ methods will instead provide a probability density function (PDF) for each source, which allows for more information on how likely the point estimate is for a given photo-$z$ algorithm. If photometric redshifts were perfectly calibrated, the convolution of these PDFs would accurately describe the PDF $n(z)$ for the tomographic bins. This is not possible in practice due to the performance limitations of photo-$z$ algorithms with limited spectroscopic training sets. 

Thus, challenges arise when debiasing and calibrating the redshift distributions obtained by the chosen photo-$z$ algorithm(s), and as a result, weak lensing cosmological analyses can suffer from this systematic. Differences in tomographic redshift calibration can lead to significant cosmological parameter shifts, as demonstrated by Figure 1 of \cite{HSCClusteringRau2023} or by the differences in calibrations on the KiDS-Legacy dataset \cite{wright2025kidslegacyClusterZ, wright2025kidslegacyCosmoShear}, compared to the analysis performed on KiDS-1000 \cite{Asgari2021KidsCosmoShear}. 
 Another approach taken by the HSC team was to parametrize photometric redshift uncertainty using a free shift parameter of the redshift first moment for the higher redshift tomographic bins lacking good redshift calibration. While this can encompass systematics, it in turn weakened the obtained cosmological information on $S_8$ and $\Omega_m$ \cite{CosmoShearLi2023HSCY3, Dalal2023CosmoShearHSC}. Moreover, this self-calibration approach requires that the parameterization of the error in the redshift distribution is accurate.

To mitigate the weaknesses of photo-$z$ algorithms and to debias the photo-$z$ distributions, a number of methods have been developed to calibrate the sample redshift distribution of each bin. 
A popular approach is to rely on the optical photometry information given by the survey and known spectroscopic redshifts from external datasets to infer redshift distributions from the galaxy population color-space, often using data reduction algorithms and machine learning methods. This technique is commonly executed with Self-Organizing Maps (SOM) \cite{WrightSOMz, SOMPZDES2024, Hildebrandt2021ClusteringRedshiftsSOM, Masters2017C3R2SOM, Buchs2019SOMRedshifts}, or through direct weighted calibration \cite{LimaDIRMethod2008}. 
The limiting factor of this approach is often the source depth and color-space coverage of spectroscopic surveys, as spectroscopic redshift targets are more prone to failures and need longer integration times on faint targets with few easily discernible features for spectroscopy. 

Another method is to use the clustering of galaxies: by leveraging angular cross-correlations between a spectroscopic sample in small redshift slices and the galaxies in each of the tomographic bins of the weak lensing survey dataset, the technique known as clustering redshifts allows to debias \cite{Newman2008ClusterZ,ménard2014clusteringbasedredshiftestimationmethod, Gatti2018DESY1ClusteringRedshifts, Gatti2018DESY1systematics,Gatti2021ClusteringDESWLsrcDistrib, Cawthon2022DESY3Boss,Davis2017DESyear1ClusterZWLsrcdistrib,EuclidClusterRedshifts2025,KidsCrossCorrvdB2020,wright2025kidslegacyClusterZ,HSCClusteringRau2023,Morrison2017TheWiZZClusterZ, Schmidt2013ClusteringRedshifts, McQuinnClusterZ2013} the $n(z)$ distributions. 
This is the method employed in this work. 
However, clustering redshifts is also subject to systematics, since the method benefits from information of small angular scales (typically $0.1-5\hMpc$) yielding higher signal to noise, though galaxy bias modeling at these scales remains difficult. 
In this work, common possible systematics that can affect and bias the clustering redshifts calibration of the photometric sample are investigated, such as the impacts of galaxy bias and magnification. 
This methodology is applied to the S19A Shape Catalog \cite{HSCShapeCataloguePDR3Li2022} from the Hyper Suprime-Cam Subaru Strategic Survey Public Data Release 3 \cite{HSC2022datarelease3} (S19A HSC-SSP PDR3) as our photometric catalog to calibrate. 
Hereafter, we refer to this dataset simply as HSC. 
The Dark Energy Spectroscopic Instrument Data Release 1 \cite{DESI.DR1.I.Presentation} and Data Release 2 \cite{DESI.DR2.Presentation.InPrep} (DESI DR1, DR2) will serve as our spectroscopic calibration samples.   

The Dark Energy Spectroscopic Instrument (DESI) is a Stage IV spectroscopic survey \cite{DESI2016b.Instr, DESI2022.KP1.Instr, Corrector.Miller.2023} carrying out a eight year program that started collecting data in 2021. DESI uses a focal plane hosting 5000 robotic positioners containing optical fibers \cite{FiberSystem.Poppett.2024} to obtain spectra, and therefore accurate spectroscopic redshifts \cite{Spectro.Pipeline.Guy.2023}, from galaxies. DESI plans to probe $17,000$ deg$^2$ of the sky in its main phase of operations \cite{SurveyOps.Schlafly.2023}, covering spectroscopic redshift ranges up to $z\sim4.2$. The DESI experiment studies six different types of tracers: the Milky Way Survey, observing stars in our local galaxy, the Bright Galaxy Survey (BGS, $0.01\lesssim z\lesssim0.6$) \cite{DESI.Selection.BGS.Hahn2023}, the Luminous Red Galaxies (LRG, $0.3\lesssim z\lesssim1.2$) \cite{DESI.Selection.LRG.Zhou2023}, the Emission Line Galaxies (ELG, $0.7\lesssim z\lesssim1.6$) \cite{DESI.Selection.ELG.Raichoor2023}, the Quasi-Stellar Objects (QSO, $0.7\lesssim z\lesssim2.1$) \cite{DESI.Selection.QSO.Chaussidon2023} and the Lyman-$\alpha$ forest (Ly$\alpha$, $z\gtrsim 2.1$). With its accurate redshift catalog, DESI has already allowed for significant breakthroughs on cosmological measurements \cite{DESI.DR2.II.BAO, DESI2024.VII.KP7B}.

The Hyper Suprime-Cam Subaru Strategic Program\footnote{\href{https://hsc-release.mtk.nao.ac.jp/doc/}{\textcolor{blue}{hsc-release.mtk.nao.ac.jp/doc/}}}\cite{HSCOverviewAihara2017} is a Stage 3 optical imaging survey using the Hyper Suprime-Cam wide-field  camera installed on the Subaru $8.2$m telescope on the summit of Maunakea, Hawai'i, and covers three depth layers: the Wide field, the Deep field, and Ultra-Deep field. The S19A (from the the internal September 2019 data release) Shape Catalog \cite{HSCShapeCataloguePDR3Li2022} compiles galaxy shapes from $i_{\mathrm{band}}$ imaging data obtained from 2014 to 2019 in the Wide field of the HSC-SSP. It is crucial to obtain well calibrated redshift distributions of this dataset, as constraints from cosmic shear heavily depend on this larger area.

In the tomographic bin analysis performed by the HSC collaboration \cite{HSCClusteringRau2023}, there was no calibration from angular clustering of galaxies above $z\gtrsim1.2$ due to the redshift limitations of the LRG dataset used in that study, therefore not providing calibration information from clustering redshifts for most of Bin 3 and none for Bin 4. 
On the color-redshift calibration aspect, lack of high redshift spectroscopic sources and deep infrared photometry led to degeneracies that made for difficult calibration.
In turn, this makes the mean value from the fiducial calibration of a given redshift bin uncertain, meaning that major shifts are possible on Bins 3 and 4.
Indeed, analysis by \cite{Zhang2022PhotometricRedshiftShifts} showed a single shift parameter was sufficient to capture $n(z)$ uncertainty per bin in the HSC distributions.
The cosmic shear analyses on the two-point correlation function \cite{CosmoShearLi2023HSCY3} and on the angular power spectrum \cite{Dalal2023CosmoShearHSC} of HSC used a conservative uniform prior $\mathcal{U}([-1,1])$ for the tomographic bins redshift shift parameters.
Both analyses report hints of photo-$z$ miscalibration, with possible redshift expectation values being higher than initially obtained from the redshift bins only calibrated with photometry. Reported marginalized mode of shifts with asymmetric $34\%$ confidence interval and maximum a posteriori (MAP) values of shifts obtained by both HSC cosmic shear analyses have returned\footnote{Note the following convention: a negative shift implies the found expectation value (from cosmic shear analysis) for a bin is higher than the available measurements.}: %$\dz=\langle z\rangle_{\mathrm{meas}} - \langle z\rangle_{\mathrm{true}}$
$\dz_3=-0.075^{+0.056}_{
-0.059} (-0.046)$ and $\dz_4=-0.157^{+0.094}_{-0.111} (-0.144)$ 
for the angular power spectrum \cite{Dalal2023CosmoShearHSC} and 
$\dz_3=-0.115^{+0.052}
_{-0.058} (-0.120)$ and $
\dz_4=-0.192^{+0.088}
_{-0.088} (-0.190)$ 
for the angular two point correlation functions \cite{CosmoShearLi2023HSCY3} in Bin 3 and Bin 4. 
More recent works \cite{Rana2025HSCY3CosmicShearRatios} used Cosmic Shear Ratios (CSR) in order to provide calibration for the photometric redshift bins, while \cite{Zhang2025RedshiftGCPointMass} leverages galaxy clustering and weak lensing (GC-WL) to infer constraints on the source redshift bins. 
While the results found by that calibration are consistent with those obtained by the cosmic shear analyses, the error bars loosely constrain the redshift expectation: the obtained results were $\Delta z_3 = -0.002^{+0.085}_{-0.217},\;\Delta z_4 = -0.292^{+0.229}_{-0.324}$ for the CSR study and $\Delta z_3 =-0.112^{+0.046}_{-0.049},\;\Delta z_4 =-0.185^{+0.071}_{-0.081}$ for the GC-WL analysis. 
In this work, the focus will be on constraining these two shift parameters, comparing results to the HSC analyses as well as previous calibrations, and producing $n(z)$ distributions for each of the bins. 

The outline of this paper is as follows. 
Section \ref{sec:data} presents the datasets employed in this study: HSC’s Shape catalog and DESI DR1, DR2, as well as the quality cuts applied to both catalogs. 
Section \ref{sec:modeling} describes the modeling of the $n(z)$ distributions with two-point correlation functions and discusses systematic effects, including galaxy bias from both samples and magnification effects. 
Section \ref{sec:2pcf} describes our measurements of the angular correlation vectors. 
Section \ref{sec:parametrization} presents a spline based approach to model the $n(z)$ distributions and discusses other methods to regularize the clustering redshifts measurements. 
Finally, results are showcased in section \ref{sec:results}.